%! Carrying Out a Test for the Difference of Two Population Means — Unit 7, Lesson 7.9 — Warm-Up (Answer Key)
\documentclass[10pt]{article}
\usepackage{apstats-article}
\usepackage{apstats-key}   % pulls in -boxes; do NOT also load -boxes
\newcommand{\TallMath}[1]{$\displaystyle #1\rule[-1.4em]{0pt}{3.2em}$}

\begin{document}

\pageheader{Unit 7, Lesson 7.9}{Warm-Up --- Answer Key}
\namedateperiod

\vspace{0.1in}

\begin{notesbox}{Spiral Review --- Setting Up a Two-Sample $t$-Test (Lesson 7.8)}

A school district wants to know whether students who attend after-school tutoring have
a higher mean math score than students who do not. A random sample of 20 tutored students
gives $\bar x_1 = 84.2$, $s_1 = 9.1$; a random sample of 22 non-tutored students gives
$\bar x_2 = 79.6$, $s_2 = 10.3$. Dotplots for both groups are roughly symmetric with no outliers.

\textbf{1.} Name the inference method.

\quad Method: \ans{Two-sample $t$-test for a difference of means}

\textbf{2.} State hypotheses. (Let Group 1 = tutored, Group 2 = non-tutored.)

\quad $H_0$: \ans{$\mu_1 - \mu_2 = 0$} \quad $H_a$: \ans{$\mu_1 - \mu_2 > 0$}

\textbf{3.} Identify the degrees of freedom range. (Use the conservative formula.)

\quad $df$ is between \ans{19} and \ans{40}

\end{notesbox}

\begin{notesbox}{Spiral Review --- Computing a One-Sample $t$-Statistic (Lesson 7.5)}

A researcher claims the mean resting heart rate of adult swimmers is 62 bpm.
A random sample of 16 swimmers gives $\bar x = 59.4$ bpm, $s = 4.8$ bpm.
All conditions are met.

\textbf{4.} Compute the $t$-statistic. Show your work.

\vspace{0.05in}
$SE = \dfrac{s}{\sqrt{n}} = \dfrac{{\color{keyred}\mathbf{4.8}}}{\sqrt{{\color{keyred}\mathbf{16}}}} = $ \ans{1.2}
\hfill
$t = \dfrac{\bar x - \mu_0}{SE} = \dfrac{{\color{keyred}\mathbf{59.4}} - {\color{keyred}\mathbf{62}}}{{\color{keyred}\mathbf{1.2}}} \approx$ \ans{$-2.17$}

\end{notesbox}

\end{document}
